\nuevaReceta{Empedrao Almeriense}{4 personas}{1 h 15 min}{receta:empedrao}

    \begin{minipage}[t]{0.35\textwidth}
        \section*{Ingredientes}
        \begin{itemize}[leftmargin=*]
            \item 250 g de arroz
            \item 300 g de habichuelas cocidas
            \item 250 g de bacalao desalado
            \item 1 Cebolla
            \item 1 Pimiento verde
            \item 2 Tomates maduros
            \item 2 Dientes de ajo
            \item 1 Pimiento seco o ñora
            \item 1 Hoja de laurel
            \item Pimentón dulce
            \item Comino
            \item Aceite de oliva
            \item Agua o caldo de pescado
            \item Sal
        \end{itemize}
    \end{minipage}
    \hfill
    \begin{minipage}[t]{0.6\textwidth}
        \section*{Preparación}
        \begin{enumerate}
            \item \textbf{Preparar el bacalao:} Desalarlo con antelación, secarlo y cortarlo en trozos medianos.
            \item \textbf{Hacer el sofrito:} Cocinar en una cazuela con aceite la cebolla, los ajos y el pimiento verde picados. Añadir el tomate rallado y dejar que reduzca bien.
            \item \textbf{Preparar el majado:} Hidratar el pimiento seco y machacar su carne con una pizca de comino. Incorporarlo al sofrito junto con el pimentón.
            \item \textbf{Formar el caldo:} Añadir las habichuelas, el laurel y el agua o caldo de pescado. Cocer a fuego suave durante 15 minutos para integrar los sabores.
            \item \textbf{Añadir el arroz:} Incorporar el arroz y cocer durante unos 10 minutos, removiendo ocasionalmente.
            \item \textbf{Incorporar el bacalao:} Añadir los trozos de bacalao y continuar la cocción otros 8-10 minutos, hasta que el arroz esté en su punto.
            \item \textbf{Reposar:} Ajustar de sal con cuidado, retirar del fuego y dejar reposar 5 minutos antes de servir.
        \end{enumerate}
    \end{minipage}

    \vspace{1em}
    \begin{tcolorbox}[colback=orange!5, colframe=orange!50, title=Consejo, size=small]
        Prueba el caldo antes de añadir sal, ya que el bacalao puede aportar suficiente. El empedrao debe quedar jugoso y ligeramente trabado por las habichuelas.
    \end{tcolorbox}
\end{receta}
